\renewcommand{\vec}[1]{\overrightarrow{\mathbf{#1}}}

\section{Fully Homomorphic Encryption}
\begin{itemize}
    \item Capabilities of (F)HE
    \item Maybe: Difference between leveled, fully, etc.
    \item We care about polynomials in the ring $R_q = \mathbb{Z}_q[X]/(X^N-1)$
    \item \gls{rlwe} decision/search problem?
    \item Focus on BGV (and TFHE in next section)
\end{itemize}

%%
%% RLWE
\begin{definition}[RLWE ciphertext] $\vec{x} \in R_q^2$ containing a (randomly generated) \textit{public element} $a$ and secret element $a\cdot s + \Delta m + \epsilon$ derived from a persistent secret key $s$ and randomly sampled noise $\epsilon$, where $\Delta$ is the noise scale and $m\in R_t$ is the plaintext message being encrypted.
    \begin{gather} % \vec{x}
        \rlwe(m) = (b, a) = (a\cdot s + \Delta m + \epsilon, a) \\
        s \overset{\$}{\leftarrow} \{-1,0,1\}^N,\quad \epsilon \overset{\$}{\leftarrow} \chi, \quad a \overset{\$}{\leftarrow} \chi
    \end{gather}

    An RLWE ciphertext $\vec{x}$ encrypting a message $m$ under the secret key $s$ is denoted $\rlwe_s(m)$, but if the secret key is understood from context then $\rlwe(m)$ is typically used.
\end{definition}

\begin{itemize}
    \item Most FHE schemes support homomorphic addition and multiplication, and admits one of these for \textit{free}, in the sense that...
    \item For example RSA has free multiplication, BGV has free addition. 
\end{itemize}

%%
%% BGV ADDITION
\begin{definition}[BGV addition.]
    \begin{equation}
        \vec{x} + \vec{y} = (b_x, a_x) + (b_y, a_y) = \rlwe_s(m_x + m_y)
    \end{equation}
    \begin{proof}
        \begin{gather*}
            % \vec{x} + \vec{y} = (b_x, a_x) + (b_y, a_y) \\
            (a_x s + \Delta m_x + \epsilon_x, a_x) + (a_y s + \Delta m_y + \epsilon_y, a_y) \\
            ((a_x +a_y )s + \Delta (m_x+m_y) + (\epsilon_x + \epsilon_y), a_x + a_y) \\
            \text{Let } a = (a_x + a_y) \text{ and } \epsilon = (\epsilon_x + \epsilon_y)\\
            \therefore \vec{x} + \vec{y} = (as + \Delta(m_x + m_y) + \epsilon, a) = \rlwe_s(m_x + m_y)
        \end{gather*}
    \end{proof}
\end{definition}

\begin{itemize}
    \item BGV multiplication is more complicated + expensive and not used in sPAR. 
    \item Requires relinearization to manage noise without bootstrapping.
\end{itemize}

%%
%% MODULUS SWITCHING
\begin{definition}[Modulus Switching]
    Explained now because its different in RNS (important)
\end{definition}

%%
%% OTHER
\begin{definition}[Boostrapping]
    Homomorphically evaluate decryption circuit to reset noise.
\end{definition}

\begin{definition}[Keyswitching]
    The process of converting a message to being encrypted under a different secret key.
\end{definition}

\newpage